\documentclass{beamer}

\usepackage{fontspec}
\newfontfamily{\monaco}{Monaco}
\newfontfamily{\comicsans}{Comic Sans MS}
\newfontfamily{\comicsansbf}{Comic Sans MS Bold}

% \pdsetup{lf={Theory and Practice of Functional Programming}}
\useinnertheme[shadow]{rounded}
\useoutertheme{infolines}

\setbeamertemplate{frametitle}{
  \comicsansbf{\insertframetitle}
}
\title[Foundations of Programming Languages]{
{\small \textcolor{black}{\textsc{Foundations of Programming Languages}}} \\[1em]
\comicsansbf{Theory and Practice of Functional Programming}
}
\author[Y. ZHANG]{\monaco{Dr. ZHANG Yu}}
\institute[ISCAS]{Institute of Software, Chinese Academy of Sciences}
  %% \\[5pt] \includegraphics[height=.04\slideheight]{iscas.eps} 
  %% \qquad \qquad 
  %% \includegraphics[height=.04\slideheight]{cas.eps}
\date[GUCAS 2011]{Fall term, 2011 \qquad GUCAS, Beijing}

\begin{document}


\begin{frame}
\titlepage
\end{frame}

%% \begin{frame}[toc=]{Outline}
%% \tableofcontents
%% \end{frame}

\begin{frame}{Theory and Practice of Functional Programming}
The course is officially entitled {\em Foundations of Programming Languages}, but I decided to focus on 
\emph{functional languages} this year --- not just as a theoretical model for PLs, but also as practical languages 
to build real applications.
\begin{itemize}
\item Functional languages are becoming popular in recent years and used to build real applications in industry:
\begin{itemize}
\item Erlang has been used for almost two decades in Ericsson; Twitter switched to scala two years ago;  
  part of Amazon Web Services is written in Erlang ...
\item Finance industry starts to make intensive use of FPL: 
  Erlang in Goldman Sachs and Deutsche Bank, Ocaml in Jane Street, Scala in UBS Investment Bank, 
  F\# and Haskell in Credit Swiss, ...
\end{itemize}
\item Programming in FPL can help to understand abstract concepts in theory.
\end{itemize}
\end{frame}

\begin{frame}{Why Haskell?}
We shall use Haskell as the main language in the course (with a comparison to Ocaml).
\begin{itemize}
\item Haskell is known to be hard for non-FPL programmers.
\item Haskell is rich, covering almost all theoretical concepts that we will introduce in the course.
\item Haskell is a pure functional language --- it is \emph{elegant}!
\end{itemize}
\end{frame}

\begin{frame}{Course outline}
10 lectures in 3 parts:
\begin{itemize}
\item \emph{Foundation and theory} of FPL.
\begin{itemize}
\item Lambda-calculus, simple type theory, semantic model, modeling of imperative programs, ...
\item Basic category concepts, monads, concurrency, ...
\end{itemize}
\item \emph{Practice programming} using FPL.
\begin{itemize}
\item Programming in Haskell: lists, type classes, module systems, functors, ...
\item Data structure and algorithms in FPL.
\end{itemize}
\end{itemize}
\end{frame}

\begin{frame}{Schedule}
\hspace*{4em}
\begin{tabular}[b]{l || l }
\hline \\[-1em]
Week 1 & Lambda calculus
\\[0.2em] \hline \\[-1em]
Week 2  & Simply typed lambda-calculi (PCF)
\\[0.2em] \hline \\[-1em]
Week 3 & Type theory 
\\[0.2em] \hline \\[-1em]
Week 4 & Introduction to Haskell 
\\[0.2em] \hline \\[-1em]
Week 5, 6 & Functional data structure and algorithms
\\[0.2em] \hline \\[-1em]
Weak 7, 8 & Advanced Haskell programming 
\\[0.2em] \hline \\[-1em]
Weak 9 & Advanced topics in type theory
\\[0.2em] \hline \\[-1em]
Week 10 & Course review
\\[0.2em] \hline 
\end{tabular}
\\
This is a rough schedule --- after the introduction of Haskell, theory and programming practice will be mixed.
\end{frame}

\begin{frame}{Reference}
\begin{itemize}
\item There is no textbook for the course. I list some reference books on the course page:
\begin{center}
\url{http://lcs.ios.ac.cn/\~yzhang/fopl/}
\end{center}
More references will be given during the course and available on the page, 
but you are also encouraged to do some research \emph{on your own}.
\\
Lecture notes and other material will  be download-able from the page.
\item Me: \\
I will be available for answering your questions {\em every Tuesday afternoon from 15:00 to 18:00}, 
at \emph{Room 315 of ISCAS building NO. 5}.
\end{itemize}
\end{frame}

\begin{frame}{Evaluation}
Evaluation will base on 
\begin{itemize}
\item {\bf Projects}: there will be small projects on functional programming. 
  You will do some implementation and make short reports. 
\item {\bf Final exam}: this consists of all the theoretical contents that are covered by the course, 
  as well as functional programming concepts. \\
  The exam will take place on \emph{Friday January 13, 2012}.
\end{itemize}
\end{frame}

\end{document}
 

